% ======================================================================================================
\section{Introduction}\label{sec:intro}
% ======================================================================================================

In an electron storage ring the fundamental-mode RF system and the stored beam form one coupled dynamical system. The beam induces voltage in the accelerating cavities; the cavities are detuned to compensate the reactive part of that loading; feedback around the klystron lowers the impedance that the beam sees; and the coherent dipole motion of the bunches responds to every perturbation of the gap voltage that the feedback leaves. The classical analyses---Robinson's stability criteria~\cite{Robinson1964}, Pedersen's beam-loaded cavity~\cite{Pedersen1975}, and the treatments of RF feedback by Boussard~\cite{Boussard1985} and Wilson~\cite{Wilson1993}---describe this coupling with a handful of parameters: cavity bandwidth and detuning, loop gain, phase and delay, beam current and synchrotron frequency. Detailed models built on the same physics were used to design and commission the low-level RF (LLRF) systems of PEP-II and the LHC~\cite{Rivetta2007,Mastorides2010}.

At an operating machine, however, several of these parameters are not known well. Components age and are replaced, calibrations are uncertain, parts of the loop setting are not recorded, and the klystron runs close to saturation, where its small-signal response to a change of drive amplitude differs from its response to a change of drive phase. A model assembled from design values may describe one operating point and fail at another. The operational questions---how the damping and frequency of the coherent mode, and the transfer of RF noise to the beam, change with beam current, loop gain and loop phase---call for a model identified from measurements. Such measurements must be made on a running machine: with the feedback closed, with stored beam, with small excitation, and preferably with the diagnostics that already exist~\cite{Forssell1999,Ljung1999}.

This paper asks whether operational closed-loop measurements can identify a compact model that predicts beam--cavity--amplifier--feedback interactions across operating conditions. The approach is physics-informed (grey-box) identification~\cite{Bohlin2006,Pintelon2012}. Physics supplies the structure---cavity resonators, rigid-bunch dipole dynamics, the loop topology and the controller networks, the steady-state beam loading, the way the loop settings enter---and the data determine, and test, the parameters that the physics leaves uncertain. The emphasis is not on fitting one transfer function, which many models do, but on whether one set of physical parameters accounts for the changes of the dynamics as the beam current, the loop gain and the loop phase change, and whether the identified model predicts quantities that were not used to identify it.

The data are the closed-loop responses recorded by the existing routine of the SPEAR3 RF station---band-limited noise injected into the klystron drive and the feedback signal recorded---at beam currents from 20 to \SI{500}{mA} and at several loop gains and phases, in two sessions a month apart, together with records that the station froze at trips of its present klystron. Four features make the problem a stringent test: (i) the excitation is one-dimensional, the same noise sequence on both quadratures, so each record measures a single projection of a response that is not rotation-invariant once the klystron is compressed; (ii) part of the loop phase is not recorded, because a tracking correction drifts by several degrees; (iii) the beam's coherent mode occupies a few kilohertz around $\pm f_s$, a band the routine itself discards; and (iv) independent information on the klystron exists only in those trip records, which are five years older than the closed-loop data.

The paper makes four contributions. Section~\ref{sec:model} gives a two-quadrature (widely linear~\cite{Picinbono1995}) model of the closed loop with beam, including the klystron's compression and AM/PM conversion, and derives what a one-dimensional excitation records. Section~\ref{sec:ident} shows that single records do not determine the klystron's asymmetry, that records taken at different loop settings determine its compression but not its AM/PM conversion, and that the trip records supply the latter. Section~\ref{sec:explain} tests competing explanations for the change of the response with beam current. Section~\ref{sec:valid} validates the selected model by out-of-band prediction of the coupled beam mode, by prediction across operating conditions, and against the klystron compression measured in the trip records; the out-of-band test exposes, and the near band then identifies, the one beam parameter that the loop band cannot determine---the synchrotron frequency.

% ======================================================================================================
\section{The measured system}\label{sec:system}
% ======================================================================================================

\subsection{RF station and direct feedback}

SPEAR3 is a \SI{3}{GeV} storage ring with harmonic number $h=372$ and $\fRF=\SI{476.3}{MHz}$, operated at \SI{500}{mA} in top-off~\cite{Hettel2002,McIntosh2004}. One \SI{1.5}{MW} klystron drives four single-cell PEP-II cavities~\cite{Rimmer1995,Rimmer2000} through a circulator and a waveguide splitter; together they provide $V_\mathrm{RF}=\SI{2.85}{MV}$. Table~\ref{tab:params} lists the parameters used. The cavities are overcoupled and are detuned below $\fRF$ in proportion to the beam current, by about \SI{36}{kHz} at \SI{500}{mA}, to cancel the reactive beam loading; the tuner loops hold the load angle, so the detuning follows the current.

Two cavity and ring constants follow from station snapshots, given the cavities' measured $R/Q=\SI{\RoverQ}{\ohm}$ and $Q_0=\Qzero$~\cite{Rimmer2000}, i.e.\ a shunt impedance $R=V^2/2P=\SI{\Rshunt}{M\ohm}$. Without beam and on tune, the reflection coefficient of each cavity is $(\beta-1)/(\beta+1)$; the snapshots give couplings $\beta=\betaLo$--\betaHi, loaded quality factors $Q_0/(1+\beta)=\QLreflLo$--\QLreflHi, and \QLrefl{} for the vector sum of the four (half-bandwidth \SI{\halfBW}{kHz}); their forward minus reflected power is \wallLo--\wallHi{} times the wall loss $V^2/2R$ computed from the probe voltages. At \SI{500}{mA} the power balance of the four cavities, $\sum_k(P_{\mathrm{fwd},k}-P_{\mathrm{refl},k}-V_k^2/2R)/\Ib$, gives the energy loss per turn $U_0=\UzeroLo$--\SI{\UzeroHi}{MeV}: the \SI{\UzeroDip}{MeV} of the dipoles~\cite{Hettel2002} plus the insertion devices. We use $U_0=\SI{\Uzero}{MeV}$, hence $\phi_s=\ang{\phiS}$ from the crest, and the radiation damping rate $\arad=J_sU_0/(2E_0T_\mathrm{rev})=\SI{\alphaRad}{s^{-1}}$ with $J_s=\Js$, the dipoles' partition number 1.79~\cite{Hettel2002} weighted with 2 for the insertion devices, which sit in dispersion-free straights~\cite{Sands1970}.

The legacy LLRF processor (RFP), an analog VXI module of PEP-II design~\cite{Corredoura1999}, regulates the vector sum of the four probe signals with a direct loop: the measured vector sum is rotated and scaled by a baseband modulator (the loop-phase and loop-gain settings; a frequency-offset tracking correction adds to the phase), the reference is subtracted, and a lead network and a lag (integral) network shape the error before the I/Q modulator that drives the klystron (Fig.~\ref{fig:loop}).

\begin{figure}[tb]
\centering
\resizebox{\columnwidth}{!}{% Block diagram of the measured closed loop with beam (column width).
\begin{tikzpicture}[font=\footnotesize, >=Latex, node distance=4mm and 5mm,
  blk/.style={draw, rounded corners=1pt, minimum height=7mm, minimum width=10mm, align=center, fill=white},
  sum/.style={draw, circle, inner sep=0pt, minimum size=4.5mm},
  sig/.style={->, thick},
  lbl/.style={font=\scriptsize, align=center}]
  % forward path
  \node[sum] (s1) {};
  \node[left=6mm of s1, lbl] (n) {noise $\vect n$\\(DACs)};
  \node[blk, right=of s1] (chain) {drive\\$c(s)$};
  \node[blk, right=of chain] (kly) {klystron\\$\mat K(\gA,\kAP)$};
  \node[blk, right=of kly] (cav) {cavities\\$\mat Z_{1..4}$};
  % beam
  \node[blk, below=7mm of cav] (beam) {beam\\$D_\mathrm{inc}^{-1}$};
  % feedback path
  \node[blk, below=7mm of kly] (vs) {vector\\sum};
  \node[blk, below=7mm of chain] (ctrl) {$G\mat R(\theta)$\\lead, lag};
  \node[coordinate, below=7mm of s1] (rec) {};
  \draw[sig] (n) -- node[above, lbl, pos=0.45]{$+$} (s1);
  \draw[sig] (s1) -- node[above, lbl]{$\vect d$} (chain);
  \draw[sig] (chain) -- (kly);
  \draw[sig] (kly) -- node[above, lbl]{$\vect i_g$} (cav);
  \draw[sig] ([xshift=2mm]cav.south) -- node[right, lbl]{$\hat{\vect p}^\mathsf{T}\!\sum\vect v_k$} ([xshift=2mm]beam.north);
  \draw[sig] ([xshift=-2mm]beam.north) -- node[left, lbl]{$I_{b1}\hat{\vect b}\Delta\phi$} ([xshift=-2mm]cav.south);
  \draw[sig] (cav.east) -- ++(3mm,0) |- ([yshift=-11mm]beam.south) -| (vs.south) node[pos=0.25, below, lbl]{probes $\vect v_k$};
  \draw[sig] (vs) -- (ctrl);
  \draw[sig] (ctrl) -- node[above, lbl]{$\vect x$} (rec) -- node[right, lbl]{$-$} (s1);
  \draw[->, dashed] (rec) -- ++(-9mm,0) node[left, lbl]{TOTAL\\recorded};
\end{tikzpicture}
}
\caption{The measured system. Band-limited noise $\vect n$ is added to the klystron drive after the controller; the controller output $\vect x$ is recorded. The beam couples to the four cavities through its dipole motion $\Delta\phi$ [Eq.~\eqref{eq:beam}]; the klystron responds differently to drive amplitude and phase [Eq.~\eqref{eq:klystron}]. Bold symbols are two-quadrature vectors and matrices in the cavity frame.}\label{fig:loop}
\end{figure} The loop delay, measured as \SI{\tauMeas}{ns}~\cite{Sebek2026}, lies mostly in the klystron and the cables to the tunnel. A digital ripple loop acting on the 60-Hz harmonics up to \SI{2.9}{kHz}, and slow EPICS loops for the klystron voltage, the cavity tuners and the gain tracking, complete the system. The klystron-voltage loop holds the drive power at a set point, \SI{\driveSet}{W}, \SI{\driveMargin}{dB} below the saturation setting of \SI{\driveSat}{W}, by moving the HVPS voltage, so the klystron's voltage, and with it its gain, change with the beam current. In the station snapshots the drive is \SI{\driveLow}{W} without beam and \SI{\driveHigh}{W} at \SI{500}{mA}, with the HVPS at \SI{\hvpsLow}{kV} and \SI{\hvpsHigh}{kV}, a klystron gain of \SI{\klyGainLow}{dB} and \SI{\klyGainHigh}{dB}, and an output of \SI{\klyFwdLow}{kW} and \SI{\klyFwdHigh}{kW} at an efficiency of \SI{\klyEffLow}{\percent} and \SI{\klyEffHigh}{\percent}: at the same drive the tube works much closer to saturation with beam than without. Gain tracking holds the product of a DSP coefficient and the measured klystron gain constant; in the same snapshots it is \gaintrackHigh{} at \SI{500}{mA} and \gaintrackLow{} without beam, although the klystron gain differs by \SI{\klyGainDiff}{dB}. The forward-path gain is therefore independent of the operating point, and the loop gain changes only with its setting. Neither the closed-loop records nor the trip records carry the HVPS voltage or the klystron's drive and output power.

\begin{table}[tb]
\caption{Parameters of the SPEAR3 RF station used in the model. Fixed values come from the design report and published measurements~\cite{Hettel2002,Huang2007,Rimmer2000,Corredoura1999} and from station snapshots (Sec.~\ref{sec:system}); identified values are those of Sec.~\ref{sec:explain}.}\label{tab:params}
\footnotesize
\begin{ruledtabular}
\begin{tabular}{lll}
Quantity & Symbol & Value \\
\hline
Energy; harmonic number & $E_0$; $h$ & \SI{3.0}{GeV}; 372 \\
RF frequency & $\fRF$ & \SI{476.304}{MHz} \\
Energy loss per turn (power balance) & $U_0$ & \SI{\Uzero}{MeV} \\
Total gap voltage & $V_\mathrm{RF}$ & \SI{2.85}{MV} \\
Synchronous phase (from crest) & $\phi_s$ & \ang{\phiS} \\
Radiation damping rate & $\arad$ & \SI{\alphaRad}{s^{-1}} \\
Momentum compaction (lattice) & $\alpha_c$ & $1.18\times10^{-3}$ \\
Cavities; $R/Q$; $Q_0$ & $n_\mathrm{cav}$ & 4; \SI{\RoverQ}{\ohm}; \Qzero \\
Lead network zero; pole & & \SI{944}{kHz}; \SI{2.39}{MHz} \\
Lag network zero; pole & & \SI{3.19}{kHz}; \SI{0.26}{kHz} \\
Drive-chain pole & $\omega_a/2\pi$ & \SI{2.5}{MHz} \\
Klystron drive (set; saturation) & $P_\mathrm{d}$ & \SI{\driveSet}{W}; \SI{\driveSat}{W} \\
Klystron AM/PM (trip records) & $\kAP$ & \kapTrip \\
\end{tabular}
\end{ruledtabular}
\end{table}

\subsection{The closed-loop measurement}

The station's closed-loop measurement routine plays a stored sequence of band-limited noise through two DACs whose output is summed into the klystron drive after the controller, and records the controller output: $2^{18}$ complex samples at eight samples per turn (\SI{10.24}{MS/s}, \SI{25.6}{ms}). The stored reference is flat to $\pm\SI{500}{kHz}$. For a complex-linear loop the transfer function from the reference to the record is the complementary sensitivity $-L/(1+L)$. The routine's own estimate discards $\pm\SI{18.8}{kHz}$ around the carrier, the band of the coupled beam mode; here the transfer function is estimated anew from the stored records over the full band (Sec.~\ref{sec:estimation}). The same sequence is stored for the I and Q DACs (correlation \noiseCorr), so the excitation is one-dimensional. The measurement is non-invasive: the injected noise is a fraction of a percent of the drive and the routine runs with stored beam.

\subsection{Data}

Eighteen closed-loop records, taken in two sessions a month apart, are used; Table~\ref{tab:groups} lists them by operating point. Loop-gain and loop-phase settings are given relative to the station's operating settings. In session I the beam current was varied from 20 to \SI{497}{mA} at a loop gain of $\gainSessionOne$~dB, and at \SI{22}{mA} the phase setting was stepped over \ang{\phaseScanLow}. In session II, at about \SI{500}{mA}, the phase setting was lowered by \phaseShiftsHigh{} degrees and the gain by up to \SI{\gainDropMax}{dB}, after one record at \SI{242}{mA}. The frequency-offset tracking correction was on and is not recorded; station snapshots show it varying between \ang{+2} and \ang{+13}.

The trip records are a second, independent data set, of the same klystron, which has been in service since September 2020. At each trip the station freezes, together, the last \SI{5.0}{ms} of the klystron drive and of one cavity probe (both I/Q) and of the klystron output amplitude, all at \SI{10.24}{MS/s}. Of the \tripN{} records of trips in 2021, \tripUsed{} hold RF in every signal: \tripFull{} at full klystron output and \tripReduced{} at \tripPower\% of that power---the output the station needs without beam (\klyFwdRatio\% in the snapshots)---with the same drive to within \tripDriveDiff\%. They carry neither the beam current, nor the loop settings, nor the HVPS voltage.

\subsection{Frequency-response estimation}\label{sec:estimation}

Each record is transformed as a whole. Its discrete Fourier transform, with a resolution of \SI{\lpmBin}{Hz}, obeys at every bin $k$
\begin{equation}\label{eq:lpm}
Y_k=G_kX_k+T_k+V_k,
\end{equation}
where $X$ is the transform of the known excitation, $G$ the frequency response, $T$ a transient that varies smoothly with $k$ (the record is not periodic, so its initial and final conditions leak into every bin) and $V$ the output noise. The local polynomial method (LPM)~\cite{Pintelon2010,Pintelon2012} fits Eq.~\eqref{eq:lpm} by least squares to \lpmBins{} adjacent bins (\SI{\lpmWidth}{Hz}), with $G$ constant and $T$ quadratic in the bin offset. This removes the transient, and neighbourhoods that do not overlap give statistically independent estimates every \SI{\lpmWidth}{Hz}. The expectation of an estimate is a weighted mean of $G$ over its neighbourhood, with weights fixed by the excitation (App.~\ref{app:estimator}); the model is compared with the same weighted mean of its own response, so the comparison is exact even where the coupled mode changes the response within one neighbourhood. The noise variance follows from the residual of a second fit in which $G$ is quadratic as well, averaged over \SI{\lpmSmooth}{kHz}. The estimates are taken in the convention of the routine, the complex conjugate of $Y/X$, which makes the loop delay positive and places the high-current detuning below $\fRF$.

Two bands are distinguished: the \emph{near band}, $|f|<\SI{\fSplit}{kHz}$, which contains the carrier and the coupled mode and which the routine discards, and the \emph{loop band}, $\SI{\fSplit}{kHz}\le|f|\le\SI{\fHi}{kHz}$, where the loop dominates. The neighbourhoods tile the spectrum outward from the carrier; the one centred on it leaves out the carrier bin, which holds the record's mean, the operating point, so the estimates cover the full band without a gap. In the near band every estimate is a data point. In the loop band, where the response is smooth, four (below \SI{100}{kHz}) or eight adjacent estimates are combined with inverse-variance weights into one data point, and the model is evaluated at their mean frequency. Each record thus gives \nCellsRecord{} complex data points over the full band, \nPoints{} real ones for the eighteen records. The real and imaginary parts of the residuals, each divided by its standard deviation, define a Gaussian likelihood, so that $\chi^2/\nu\approx1$ indicates a fit at the noise level and differences of $\chi^2$ are log-likelihood ratios. The scatter of adjacent estimates confirms the predicted variance in the loop band (ratio \calibMedian, App.~\ref{app:estimator}); in the near band, interference lines raise the averaged noise level of the estimates around them, and the ratio lies between \calibNearLo{} and \calibNearHi.

The near band is much noisier than the loop band, and the records, not the estimator, set this. The excitation is flat, but within \SI{6}{kHz} of the carrier the part of the record that it does not explain---the loop's correction of other disturbances, among which the harmonics of the ripple of the klystron's high-voltage supply stand out, and, at high current, the beam's own motion---is \noiseRiseLo--\SI{\noiseRiseHi}{dB} stronger than in the loop band, and a \SI{25.6}{ms} record holds only \lpmBins{} bins per \SI{\lpmWidth}{Hz}. A single near-band estimate therefore has a median signal-to-noise ratio of \snrNearLo--\SI{\snrNearHi}{dB} (range over the records), against \snrLoopLo--\SI{\snrLoopHi}{dB} below \SI{150}{kHz} in the loop band, where several estimates are combined besides (Fig.~\ref{fig:data}). A Welch estimate with the routine's 8192-sample segments averages over about five times the bandwidth and is correspondingly smoother, but it does not resolve the coupled mode, which at high current is a few hundred hertz wide.

% ======================================================================================================
\section{Physics-informed model}\label{sec:model}
% ======================================================================================================

\subsection{Two-quadrature representation}

Signals are complex envelopes on the carrier $e^{j\wRF t}$, and $s$ is the Laplace variable of the envelope ($s=j2\pi f$ at an offset $f$ from $\fRF$). Because the klystron and the beam treat amplitude and phase differently, each envelope is written as the real vector $\vect x=(x_I,x_Q)^\mathsf{T}$ in the \emph{cavity frame}, in which the steady total gap voltage is real, and each complex-coefficient transfer function $F(s)$ as the real matrix
\begin{equation}\label{eq:realmatrix}
\mat F(s)=\begin{pmatrix}F_r & -F_i\\ F_i & F_r\end{pmatrix},\quad
F_{r,i}=\tfrac{1}{2}\left[F(s)\pm\bar F(s)\right]\big/\{1,j\},
\end{equation}
with $\bar F(s)=\overline{F(\bar s)}$, the conjugate system~\cite{Schreier2010}. Blocks with a real impulse response---delays, the lead and lag networks, the drive-chain pole---are scalars times the identity; a detuned cavity, a rotation, the compressed klystron and the beam are not.

\subsection{Cavities and beam}

Cavity $k$ is the baseband resonator $Z_k(s)=R_L\omega_{1/2}/(s+\omega_{1/2}-j\Delta\omega_k)$, with $R_L=(R/Q)\QL$, $\omega_{1/2}=\wRF/2\QL$ and $\Delta\omega_k$ its detuning from $\fRF$. All four receive the same generator current $\vect i_g$ and the same beam current, so $\vect v_k=\mat Z_k(\vect i_g-I_{b1}\hat{\vect b}\,\Delta\phi)$, where $I_{b1}=2\Ib$ and $\Delta\phi$ is the phase of the bunch centroids (positive when late). For rigid bunches in the uniform-fill dipole mode~\cite{Pedersen1975,Chao1993},
\begin{equation}\label{eq:beam}
\left(s^2+2\arad s+\omega_s^2\right)\Delta\phi=\kappa_1\,\hat{\vect p}^\mathsf{T}\textstyle\sum_k\vect v_k,
\end{equation}
with $\kappa_1=\alpha_c\wRF/[(E_0/e)T_\mathrm{rev}]$, $\omega_s^2=\kappa_1V_\mathrm{RF}\sin\phi_s$, and the unit vectors $\hat{\vect b}=(-\sin\phi_s,-\cos\phi_s)^\mathsf{T}$, along which a phase shift changes the beam current, and $\hat{\vect p}=(\cos\phi_s,-\sin\phi_s)^\mathsf{T}$, along which the bunch senses the voltage. Eliminating $\Delta\phi$ gives the beam-loaded plant from generator current to the measured vector sum $\vect v_m=\sum_kw_k\vect v_k$ ($w_k=1/4$):
\begin{align}
\mat P_b&=\mat Z_m-\frac{\kappa_1I_{b1}}{D(s)}\,\mat Z_m\hat{\vect b}\,\hat{\vect p}^\mathsf{T}\mat Z_\Sigma,\label{eq:plant}\\
D(s)&=s^2+2\arad s+\omega_s^2+\kappa_1I_{b1}\,\hat{\vect p}^\mathsf{T}\mat Z_\Sigma\hat{\vect b},\label{eq:Dopen}
\end{align}
where $\mat Z_\Sigma=\sum_k\mat Z_k$ and $\mat Z_m=\sum_kw_k\mat Z_k$. $D(s)=0$ is Pedersen's quartic for the beam-loaded cavity without feedback~\cite{Pedersen1975}.

\subsection{Klystron}

About its operating point the klystron is memoryless at these frequencies (its bandwidth is \SI{5}{MHz}). A change of drive amplitude changes the output amplitude with a relative gain $\gA$ and the output phase by $\kAP$ radians per unit relative amplitude (AM/PM conversion~\cite{Saleh1981,Cripps2006}); a change of drive phase passes with unit relative gain. In the cavity frame
\begin{equation}\label{eq:klystron}
\mat K=\mat R(\theta_g)\begin{pmatrix}\gA&0\\ \kAP&1\end{pmatrix}\mat R(-\theta_g),
\end{equation}
where $\mat R$ is a rotation and $\theta_g$ the angle of the steady generator current, computed from the beam loading [$\vect i_{g0}=V_\mathrm{cav}/Z_m(0)+I_{b1}e^{-j\phi_s}$]. A linear klystron has $\gA=1$, $\kAP=0$; a klystron driven close to saturation has $\gA<1$. Its phase-path gain is absorbed in the loop gain. Because the station holds the drive and moves the cathode voltage (Sec.~\ref{sec:system}), the klystron's gain and its place on its transfer curve change with the beam current. Gain tracking removes the change of gain from the loop; the change of place, and with it of $\gA$ and $\kAP$, remains.

\subsection{Loop and recorded response}

The noise $\vect n$ enters the drive after the controller, and the record is the controller output $\vect x$:
\begin{equation}\label{eq:loop}
\vect x=-\mat M(\vect x+\vect n),\qquad \mat M(s)=G\,c(s)\,\mat R(\theta)\,\mat P_b(s)\,\mat K/R_L,
\end{equation}
where $G$ is the loop gain of the phase path at an on-tune resonance, $\theta$ the loop phase error (positive = lead), and
\begin{equation}\label{eq:chain}
c(s)=e^{-s\tau}\,\frac{1+s/\omega_z}{1+s/\omega_p}\,\frac{s+\omega_{iz}}{s+\omega_{ip}}\,\frac{\omega_a}{s+\omega_a}
\end{equation}
collects the delay $\tau$, the lead and lag networks and the drive-chain pole (Table~\ref{tab:params}). The closed-loop response is $\mat T=-(\mat 1+\mat M)^{-1}\mat M$. A one-dimensional excitation along the angle $\alpha$ in the cavity frame records
\begin{equation}\label{eq:recorded}
H(f)=e^{-j\alpha}\,(1\;\;j)\,\mat T\begin{pmatrix}\cos\alpha\\ \sin\alpha\end{pmatrix}=H_1(f)+H_2(f)\,e^{-2j\alpha},
\end{equation}
with the direct and image responses $H_1=[T_{11}+T_{22}+j(T_{21}-T_{12})]/2$ and $H_2=[T_{11}-T_{22}+j(T_{21}+T_{12})]/2$~\cite{Picinbono1995,Schreier2010}. For a linear klystron and no beam, $H_2=0$ and $H=-L/(1+L)$ with $L=Ge^{j\theta}c\,Z_m/R_L$. Otherwise a single record measures one mixture of $H_1$ and $H_2$, set by $\alpha$, which depends on the klystron's insertion phase but not on the loop settings, because the noise enters after the controller's rotation. Elements outside the loop (the DAC filters and the recording channel) multiply $H$ by a smooth factor common to all records, $\exp(c_1x+c_2x^2)$ with $x=f/\SI{400}{kHz}$; its linear phase is not separable from the loop delay and is set to zero.

\subsection{The coupled mode}

The poles of $\mat T$ are the zeros of $\det[\mat 1+\mat M(s)]\,D(s)$. Tracked from zero current, one pair is the coherent dipole mode of beam, cavities, klystron and loop, $s=-\sigma\pm j\omega_c$. Without feedback it is Pedersen's root; with a linear klystron it is the root of the characteristic equation of the beam with an RF feedback loop~\cite{Boussard1985,Wilson1993}; App.~\ref{app:derivation} shows the exact factorization that connects the two-quadrature form to these, and the model reproduces both limits to machine precision. The same model therefore predicts the recorded response in the loop band, where the loop dominates, and in the near band, where the coupled mode does.

\subsection{Parameter classes}

The physics fixes the structure and the parameters of Table~\ref{tab:params}. The beam current of each record is measured; the loop gain is $G=G_0\,10^{\Delta G/20}$, with $\Delta G$ the gain setting in dB relative to the operating setting and $G_0$ the loop gain at the operating setting, because gain tracking holds the forward path constant. The remaining parameters are identified, at the level at which physics says they can vary:
\begin{itemize}
\item station constants, common to all records: $G_0$, $\tau$, $\QL$, the measurement-chain coefficients and, where stated, $\alpha_c$;
\item per operating point (session and current): the compression $\gA$ and the noise angle $\alpha$; the AM/PM conversion $\kAP$, which the closed-loop records do not determine (Sec.~\ref{sec:ident}), is fixed at the value measured in the trip records, and the alternative of one compression for all currents is tested in Sec.~\ref{sec:explain};
\item per group of records taken minutes apart at one current: the detuning $\Delta f$;
\item per record: the loop phase $\theta$, which carries the unrecorded tracking correction.
\end{itemize}
The steady generator-current angle $\theta_g$ follows from the detuning and the current.
